\renewcommand{\chaptertitle}{Foreword}

%  Make first page footer:
\fancypagestyle{firststyle}{%
\fancyhf{}% Clear header/footer
\fancyhead{}
\fancyfoot[L]{{\footnotesize
              %% We toggle between these:
              % Manuscript submitted for review.\\
              \authorname{}~\authorpatp{}.  ``\chaptertitle{}''.  {\it Nockonomicon} (2025):  40–$x$.
              }}
}
%  Arrange subsequent pages:
\fancyhf{}
\fancyhead[LE]{{\urbitfont Nockonomicon}}
\fancyhead[RO]{\chaptertitle{}}
\fancyfoot[LE,RO]{\thepage}

\chapter*{\chaptertitle}

\thispagestyle{firststyle}

There are many worthy raconteurs of the history of computing, such as \citet{Hofstadter1980} and \citet{Penrose2004}.\footnote{Yes, \citeauthor{Penrose2004} is talking about computing although it may not yet be what you intend by the term.}  What distinguishes my approach to the history of computing from other tellings?  Computing has deep connexions to alchemy and true language, which I hope to demonstrate without assuming any Hermetic background on the reader's part.  Much like technology, computing is not merely a tool nor a stance towards the world:  it is innate.  Wolfram presented this as his Principle of Computational Equivalence, ``Almost all processes that are not obviously simple can be viewed as computations of equivalent sophistication.''\footnote{\citet[p.\ 5]{Wolfram2002}.}  However, this stance does not entail technologism's supercession; computing is therefore always both a lens and a kaleidoscope, through which many layers of reality can be seen simultaneously.  It encounters the world as a system of pure statements---or purifying statements, statements hauling themselves via adjutant humanity out of incoherence into nitid clarity---and acting to make those statements true.  Mankind as unwitting amanuensis of the immanent.  Apocalypse---unveiling---as both ontology and epistemology.

First, therefore, I must establish that computation is the apex of a long human search for a language of truth.  The origin of symbol itself in the σύμβολον, a token broken in half to infallibly mark identity, is one font.  Later ``symbol'' was ``credo,'' the profession of faith, both an attempt to formulate truth and a sign of shared claims.  The primary source, however, is language itself, Hofstadter's ``strange loop'' of self-reference.  The first possible Nock formula is a crash, the void.  The second is self-reference, awakening.  Peirce would be proud of the progression.  We have tricked ourselves into thinking that computing is a late development, a fruitful elaboration of the end of mathematics.  In fact, it is the philosopher's stone, and it underlies every process in the world.  It has turned lead into gold---literally, in the case of the neutron fluences of the Manhattan Project and its battery of human computers before the advent of the vacuum tube.

The historical account which follows aims throughout at the goal of understanding what Nock is and why it matters as a protocol.  The reader will understand what mankind has been looking for, when we have caught various pieces of the puzzle, how these things came together in the mid-20th century, and how they have been refined into a pure machine.
